% =============================================================================
%  SPIE A4 Proceedings Format
%  Traffic Signal Timing Optimization Using Genetic Algorithms
% =============================================================================
\documentclass[a4paper,10pt]{article}

% ---------------------------------------------------------------------------
% SPIE A4 Margins: Top 2.54cm, Bottom 4.94cm, Left/Right 1.93cm
% ---------------------------------------------------------------------------
\usepackage[top=2.54cm,bottom=4.94cm,left=1.93cm,right=1.93cm]{geometry}

% ---------------------------------------------------------------------------
% Times Roman (matches SPIE template)
% ---------------------------------------------------------------------------
\usepackage{newtxtext}
\usepackage{newtxmath}

% ---------------------------------------------------------------------------
% Typography and layout
% ---------------------------------------------------------------------------
\usepackage{microtype}
\usepackage{setspace}
\usepackage{indentfirst}
\setlength{\parindent}{1em}
\setlength{\parskip}{0pt}

% ---------------------------------------------------------------------------
% Section heading formatting (SPIE style)
%   Section:    11pt BOLD CENTERED UPPERCASE  (e.g. "1. INTRODUCTION")
%   Subsection: 10pt bold left-justified       (e.g. "2.1 Traditional Methods")
% ---------------------------------------------------------------------------
\usepackage{titlesec}
\usepackage{textcase}

\titleformat{\section}[block]
{\centering\fontsize{11}{13}\selectfont\bfseries}
{\thesection.}{0.5em}{\MakeUppercase}

\titleformat{\subsection}[block]
{\fontsize{10}{12}\selectfont\bfseries}
{\thesubsection}{0.5em}{}

\titlespacing*{\section}{0pt}{14pt}{6pt}
\titlespacing*{\subsection}{0pt}{10pt}{4pt}

% ---------------------------------------------------------------------------
% References: [N] bracket format
% ---------------------------------------------------------------------------
\usepackage{cite}
\renewcommand{\refname}{REFERENCES}

% ---------------------------------------------------------------------------
% Hyperlinks (invisible in print)
% ---------------------------------------------------------------------------
\usepackage[hidelinks]{hyperref}

% =============================================================================
\begin{document}
	% =============================================================================
	
	% ---------------------------------------------------------------------------
	% TITLE BLOCK
	% ---------------------------------------------------------------------------
	\begin{center}
		
		{\fontsize{16}{19}\selectfont\bfseries
			Genetic Algorithm Optimization of Traffic Signals: A SUMO Microsimulation\\
			Study of Urban Intersections in Davao City, Philippines
			\par}
		
		\vspace{10pt}
		
		{\fontsize{12}{14}\selectfont
			James O. Ga-as$^{a}$, Julse M. Merencillo$^{b}$, Maureen M. Villamor$^{c}$
			\par}
		
		\vspace{4pt}
		
		{\normalsize
			$^{a,b,c}$University of Southeastern Philippines, College of Information and Computing,\\
			Obrero, Davao City 8000, Philippines
			\par}
		
	\end{center}
	
	\vspace{8pt}
	
	% ---------------------------------------------------------------------------
	% ABSTRACT
	% ---------------------------------------------------------------------------
	\begin{center}
		{\fontsize{11}{13}\selectfont\bfseries ABSTRACT\par}
	\end{center}
	
	\noindent
	Urban traffic congestion imposes disproportionate costs on rapidly urbanizing cities in Southeast Asia, where heterogeneous vehicle compositions and fixed-time signal control compound intersection delay. This study applies a genetic algorithm (GA) coupled with SUMO (Simulation of Urban MObility) microsimulation to optimize signal timing at two Davao City intersections --- E.~Quirino Avenue and the Quimpo Boulevard--GE Torres junction --- using empirical five-day DPWH traffic classification data (April 2025) as simulation demand input. Chromosomes encode green-phase durations; fitness is evaluated as $F = 1/(\sum d_i + \omega \cdot Q_{max} + \varepsilon)$ over 24-hour simulation episodes run via libsumo. \textbf{[RESULTS: Replace this sentence with your key quantitative findings, e.g., ``GA-optimized plans reduced mean vehicle delay by X\% and maximum queue length by Y\% at E.~Quirino Avenue, and by A\% and B\% at Quimpo--GE Torres, with improvements confirmed as statistically significant ($p < 0.05$, two-tailed $t$-test, $n = 10$).'']}\@ \textbf{[CONCLUSION: Replace this sentence with your main takeaway, e.g., ``Results demonstrate that simulation-coupled GA optimization is a tractable and locally adaptable approach for signal timing improvement in Philippine urban intersections with heterogeneous vehicle compositions.'']}
	
	\vspace{6pt}
	
	\noindent\textbf{Keywords:} genetic algorithm, traffic signal optimization, SUMO microsimulation, urban intersections, Davao City, Philippines, vehicle queue minimization, signal timing
	
	% ---------------------------------------------------------------------------
	% SECTION 1: INTRODUCTION
	% ---------------------------------------------------------------------------
	\section{Introduction}
	
	Urban traffic congestion imposes substantial economic and social costs on rapidly urbanizing cities in Southeast Asia. As vehicle registration rates outpace road network expansion in the Philippines, optimization of existing signal infrastructure has emerged as a cost-effective lever for improving mobility. The 2024 TomTom Traffic Index ranks Metro Manila among Asia's most congested cities, and secondary centers such as Davao City exhibit similar peak-hour gridlock [1]. Motorcycles, jeepney-type public utility vehicles, and goods trucks create a heterogeneous traffic stream whose wide variation in vehicle dynamics and saturation flow rates strains conventional fixed-time signal models [2].
	
	Fixed-time signal control dominates Philippine urban intersections: pre-programmed cycle lengths are calibrated from historical volume data and deployed without real-time adaptation [3]. More sophisticated adaptive systems such as SCOOT and SCATS have demonstrated improvements on coordinated arterials but require dense detector infrastructure and expert calibration that remain impractical for many developing-country municipalities [4].
	
	Genetic Algorithms (GAs) offer a tractable alternative for offline signal timing optimization. GAs search the feasible timing space through selection, crossover, and mutation, evaluating candidates via simulation rather than requiring a closed-form objective function [5] --- an important advantage when the objective landscape is nonlinear or noisy. Coupling a GA with SUMO (Simulation of Urban MObility) microsimulation yields timing plans that account for heterogeneous vehicle dynamics using empirical demand data as input [6].
	
	Despite a growing GA signal-optimization literature, applications using locally collected DPWH classification data at Philippine intersections remain scarce. Analytical tools calibrated for passenger-car-equivalent assumptions may underestimate the impact of the Philippines' distinctive vehicle mix on saturation flow and phase capacity [7].

	Davao City was selected as the study site because it is the largest city in Mindanao and among the fastest-growing secondary urban centers in the Philippines, with documented signal timing inflexibility at major arterials identified as a principal driver of intersection delay [7]. DPWH traffic classification data at key corridor sites in the city is available, enabling empirical demand modeling without reliance on synthetic profiles, and the existing fixed-time plans at the study intersections have not been subjected to any documented microsimulation review.

	This study addresses that gap by applying a GA-based optimizer to two Davao City intersections --- E.~Quirino Avenue and the Quimpo Boulevard--GE Torres junction --- whose existing fixed-time plans have not been subjected to any documented microsimulation review. Demand input is derived from five-day DPWH traffic classification surveys conducted in April 2025; all optimization and evaluation is conducted in SUMO models calibrated to the physical geometry, phase sequences, and vehicle mix of both sites. The specific objectives are as follows:
	
	\begin{enumerate}
		\item To construct and calibrate SUMO microsimulation models of two Davao City intersections using empirical DPWH traffic volume and vehicle classification data.
		\item To formulate signal timing optimization as a genetic algorithm problem, encoding phase green durations and cycle length as chromosomes and simulated queue length and delay as fitness criteria.
		\item To evolve GA-optimized signal timing plans for each intersection and characterize convergence behavior across generations.
		\item To compare the performance of GA-optimized plans against the current fixed-time baseline using efficiency metrics derived from full-day simulation runs.
	\end{enumerate}

	The scope of this study is limited to the two identified signalized intersections under isolated intersection control. Arterial coordination across multiple signals, pedestrian phases, non-motorized modes, and real-time adaptive deployment are outside the scope of this work.

	% ---------------------------------------------------------------------------
	% SECTION 2: REVIEW OF RELATED LITERATURE
	% ---------------------------------------------------------------------------
	\section{Review of Related Literature}
	
	\subsection{Traditional Traffic Signal Control Methods}
	
	Traffic signal control has evolved steadily since the adoption of fixed-time systems in the early twentieth century. The foundational contribution of Webster [8] established analytical formulas for computing optimal cycle lengths and green time splits as functions of approach volumes and saturation flow rates. Webster's formula, derived under the assumption of stationary Poisson arrival processes, yields deterministic phase plans that minimize average vehicle delay at isolated intersections under steady demand. This approach remains embedded in widely used signal timing tools and serves as a standard baseline against which adaptive methods are compared.
	
	Vehicle-actuated and semi-actuated controllers addressed some limitations of fixed-time systems by incorporating inductive loop detector inputs to extend or truncate green phases according to real-time presence information. Fully actuated systems vary phase durations within prescribed minimum and maximum bounds, offering improved responsiveness to short-term demand fluctuations. However, these systems optimize locally within individual phase cycles and do not learn or adapt to longer-horizon traffic patterns [3]. The absence of memory or cross-cycle coordination limits their effectiveness under recurring peak-period conditions.
	
	More sophisticated area-wide adaptive control systems, including SCOOT, SCATS, and RHODES, use real-time detector data to continuously update offsets, splits, and cycle lengths across coordinated signal networks. Srinivasan et al.~[4] reviewed neural network and adaptive approaches to signal control and found that these systems produced consistent improvements over fixed-time operation on coordinated arterials, particularly under variable demand. Nevertheless, their deployment requires dense detector infrastructure, centralized control hardware, and trained operators, creating significant barriers to adoption in developing-country municipalities. At the level of individual isolated intersections, which constitute the focus of this study, these systems do not provide advantages over simpler optimization methods when detector infrastructure is unavailable.

	\subsection{Genetic Algorithms for Signal Timing Optimization}
	
	Genetic Algorithms, introduced by Holland [5] and formalized for engineering optimization by Goldberg [9], are population-based metaheuristic search methods that represent candidate solutions as chromosomes and improve a population iteratively through selection, crossover, and mutation. Each chromosome is evaluated by a fitness function, in signal timing applications typically a simulation-derived measure of delay, queue length, or throughput, and higher-fitness individuals are preferentially selected as parents for subsequent generations. GAs are well suited to problems with non-differentiable, multi-modal, or noisy objective landscapes, characteristics routinely encountered in traffic signal optimization due to vehicle arrival stochasticity and nonlinear capacity interactions.
	
	Early applications of GAs to traffic signal timing treated phase green durations, cycle lengths, and arterial offsets as real-valued genes. Ceylan and Bell [6] applied a GA to optimize signal timings in a multi-path traffic network using TRANSYT-based simulation for fitness evaluation, demonstrating convergence to near-optimal solutions and outperforming gradient-based methods on non-convex instances. Their formulation, encoding splits and offsets as real-valued genes with Gaussian mutation, established a methodology that has remained influential in offline signal timing research.
	
	Park et al.~[10] proposed a GA-based optimization framework for isolated oversaturated intersections, evaluating fitness through microscopic simulation across multiple demand scenarios. Their results showed that GA-optimized plans reduced average delay by 12--25\% compared to Webster's analytical method under peak-hour conditions, with the stochastic search character of GAs enabling escape from local optima that gradient descent could not avoid. Importantly, Park et al.~also found that GA performance was robust to moderate demand uncertainty, suggesting that plans evolved on representative simulation demand profiles transfer reasonably well to real-world deployment.
	
	Stevanovic et al.~[11] extended GA-based optimization to coordinated arterial signal control, demonstrating that GA-optimized offset and split plans reduced stops and total travel time across arterial corridors, particularly during off-peak periods when demand patterns deviate from the historical averages used to calibrate fixed plans. Their integration of GA search with a calibrated VISSIM microsimulation model for fitness evaluation is methodologically analogous to the SUMO-coupled approach adopted in the present study.
	
	More recently, Alegre et al.~[12] demonstrated memory-assisted GA approaches for signal timing in traffic networks, encoding phase duration sequences as chromosomes and evaluating fitness using SUMO microsimulation. The use of evolutionary optimization as an outer loop over a high-fidelity simulator allows the GA to leverage realistic traffic dynamics, including heterogeneous vehicle performance, stochastic arrivals, and interaction effects between phases, without requiring explicit gradient information or a differentiable traffic flow model. This architecture is particularly appropriate for intersections with heterogeneous vehicle compositions, such as those studied here, where the nonlinear effects of motorcycle and heavy vehicle proportions on saturation flow create objective landscapes that analytical methods approximate poorly.

	While reinforcement learning agents and fuzzy logic controllers have also appeared in recent signal control literature, both require either extensive training infrastructure (RL) or expert-crafted membership functions that do not generalize across intersections (fuzzy). GA-based offline optimization, by contrast, requires only a calibrated simulation model and historical demand data, and produces a fixed timing plan directly programmable on existing signal hardware --- making it the most operationally practical choice for sites where adaptive infrastructure is unavailable.

	\subsection{Microsimulation with SUMO}
	
	SUMO (Simulation of Urban MObility) is an open-source, microscopic, and continuous traffic simulation suite developed by the German Aerospace Center (DLR). First introduced by Krajzewicz et al.~[13] and substantially extended in subsequent releases [14], SUMO supports multi-modal simulation of road networks imported from OpenStreetMap or constructed manually, with support for the Krauss stochastic car-following model, lane-changing behavior, signal logic, and diverse vehicle types parameterized by length, maximum speed, acceleration, and deceleration. Its TraCI (Traffic Control Interface) API and compiled libsumo library enable real-time and in-process interaction with running simulations, making SUMO a widely adopted platform for algorithmic traffic control research.
	
	Lopez et al.~[15] described SUMO's microscopic simulation architecture in the context of intelligent transportation research, highlighting its capacity for large-scale network simulation and suitability as a reproducible benchmark environment for traffic control algorithms. The ability to import road geometry from OpenStreetMap, define arbitrary vehicle type distributions, and script signal phase logic via XML configuration makes SUMO particularly suitable for constructing calibrated models of real-world intersections from publicly available and empirically collected data sources.
	
	The SUMO-RL library, developed by Alegre [16], wraps SUMO within the OpenAI Gymnasium interface, providing standard \texttt{reset()}, \texttt{step()}, and observation/action space abstractions compatible with optimization and learning frameworks. SUMO-RL exposes configurable observation functions encoding phase identity, per-lane queue lengths, vehicle density, and accumulated waiting times, as well as multiple built-in reward functions including queue minimization and pressure-based formulations. In the present study, SUMO-RL provided the foundational simulation infrastructure for earlier exploratory work; the GA fitness evaluation pipeline described in Section~3 interfaces with SUMO directly via the libsumo in-process library rather than through the Gymnasium action/observation abstraction, enabling full-resolution tripinfo metric extraction without the overhead of the RL step interface.
	
	A critical performance consideration in coupling GA optimization with SUMO is the computational cost of fitness evaluation: each simulation episode represents 24~hours of traffic (86,400~s) at a 15~s control step, and population sizes of 50--100 chromosomes evaluated over 30--50 generations require thousands of complete simulation runs. The use of libsumo, a compiled, in-process alternative to the socket-based TraCI interface, eliminates inter-process communication overhead and reduces simulation step time by approximately an order of magnitude relative to standard TraCI, making full-day multi-generation GA experiments computationally tractable on commodity hardware.
	
	\subsection{Traffic Signal Studies in the Philippine Urban Context}
	
	As vehicle registration growth outpaces infrastructure expansion in Philippine cities, secondary centers such as Davao City face compounding signal management challenges compounded by the absence of mass rapid transit alternatives.
	
	The Department of Public Works and Highways conducts systematic traffic classification surveys at key corridor locations throughout the Philippines, recording hourly volumes across 15 vehicle categories from motorcycles and passenger cars to articulated trucks and trailers. Data collected at E.~Quirino Avenue (Site SV11909MN\_MC, northbound, Lane 1) over five consecutive days in April 2025 reveals daily volumes ranging from approximately 18,500 to 23,300 vehicles, with consistent AM peaks of 1,220--1,544 vehicles per hour between 07:00--09:00 and broad PM peaks of 1,184--1,235 vehicles per hour between 13:00--17:00. The vehicle composition across both study intersections is dominated by passenger cars (approximately 57\%) and motorcycles (approximately 20\%), with public utility vehicles (jeepneys, 10\%), goods trucks (11\%), and buses and trailers (combined 2\%) comprising the remainder [2]. This composition differs substantially from the passenger-car-dominant distributions assumed in standard saturation flow models calibrated for North American or European traffic conditions.
	
	Reyes and Guia [7] conducted an exploratory analysis of traffic management conditions in Davao City, identifying signal timing inflexibility and the absence of demand-responsive control as principal contributors to intersection delay at major arterials. The study noted that heterogeneous vehicle composition, characteristic of Philippine urban traffic, creates challenges for signal timing models calibrated to passenger car equivalents alone, and recommended simulation-based optimization as a methodologically appropriate response to this gap. Their findings directly motivate the choice of microsimulation-coupled GA optimization in the present work, as this approach avoids the car-equivalency simplification by modeling each vehicle type with empirically grounded kinematic parameters.
	
	\subsection{Summary and Research Gap}
	
	The literature reviewed above establishes that GA-based methods consistently outperform Webster's analytical formula under peak and variable demand conditions and are robust to demand uncertainty. SUMO provides an open-source, reproducible platform that supports the vehicle type diversity required to model Philippine heterogeneous traffic accurately.
	
	Nevertheless, a specific gap exists in the application of simulation-coupled GA optimization to signalized intersections in Philippine cities using locally collected DPWH classification data. Existing international GA studies have been conducted primarily on networks calibrated for North American or European vehicle compositions and do not account for the distinctive kinematic mix of motorcycles, jeepneys, and heavy trucks that characterizes Philippine urban traffic. Furthermore, the intersections studied in existing work are typically modeled with stylized demand profiles rather than real five-day traffic counts, limiting the fidelity with which optimization results can be expected to transfer to actual deployment.
	
	This study directly addresses this gap by constructing SUMO models of E.~Quirino Avenue and the Quimpo Boulevard--GE Torres intersection from DPWH empirical data and evolving GA-optimized signal timing plans evaluated against simulated queue length and delay under realistic heterogeneous demand. The contribution is both methodological, demonstrating a complete data-to-optimization pipeline adaptable to any DPWH-surveyed intersection, and substantive, providing optimized timing recommendations for two Davao City intersections whose current fixed-time plans have not been subjected to simulation-based review.
	
	% ---------------------------------------------------------------------------
	% SECTION 3: METHODOLOGY
	% ---------------------------------------------------------------------------
	\section{Methodology}
	
	\subsection{Site Selection and Data Processing}
	Two intersections in Davao City were modeled: E.~Quirino Avenue and the Quimpo Blvd--GE Torres junction. These two intersections were selected for three reasons. First, both are identified as high-volume arterial corridors in Davao City with documented signal timing inflexibility as a principal contributor to intersection delay [7]. Second, five-day DPWH traffic classification count data is available at both sites, enabling empirical demand modeling without reliance on synthetic traffic profiles. Third, the existing fixed-time plans at both intersections --- a 150~s cycle at E.~Quirino Avenue and a multi-phase 250~s cycle at Quimpo--GE Torres --- were established without any documented microsimulation review, making them strong candidates for simulation-coupled optimization.

	We utilized 5-day traffic classification counts provided by DPWH, collected April 15--20, 2025 [2]. Vehicles were classified into five types: passenger cars (57\%), motorcycles (20\%), public utility vehicles/jeepneys (10\%), trucks (11\%), and buses (2\%). Hourly volumes were disaggregated into 15-minute arrival bins and converted to SUMO route files to capture peak-hour demand surges.
	
	\subsection{Simulation Calibration}
	SUMO network models were constructed by importing OpenStreetMap geometries into \texttt{netedit} and manually refining lane counts, channelization, and turning restrictions based on site observation records. The E.~Quirino Avenue model encodes a four-approach intersection with four controlled traffic signals; the Quimpo--GE Torres model encodes a multi-approach intersection with one coordinated traffic signal managing ten phases. Signal phase sequences were programmed to replicate the current fixed-time plans: a four-phase 150~s cycle at E.~Quirino Avenue (major green 95~s, minor green 45~s, with 5~s yellow inter-green intervals) and the existing multi-phase 250~s cycle at Quimpo--GE Torres.

	Hourly origin-destination demand was derived from the DPWH April 2025 five-day classification counts [2], disaggregated into 15-minute arrival bins and converted to SUMO \texttt{.rou.xml} route files. Vehicle type proportions (cars 57\%, motorcycles 20\%, PUVs 10\%, trucks 11\%, buses 2\%) were applied as a stochastic type distribution using SUMO's \texttt{vTypeDistribution} facility.

	To account for local driving conditions, Krauss car-following model parameters were calibrated per vehicle class. Motorcycle parameters (minimum gap $= 0.5$~m, acceleration $= 3.5$~m/s$^2$) reflect the characteristically short headways and high acceleration of Philippine motorcycle traffic. PUV/jeepney parameters were assigned reduced maximum speed (22.2~m/s) and moderate deceleration to reflect the frequent-stop dwell behavior of public utility vehicles operating along fixed routes. Heavy vehicle and truck parameters followed SUMO defaults for the goods vehicle class.

	\textbf{Baseline Validation.} Prior to GA optimization, the simulated fixed-time baseline was validated by comparing simulated queue lengths against field-observed peak-hour measurements at each intersection. Five independent baseline simulation trials were run using different random seeds, and mean per-approach queue lengths were compared against manually recorded counts. Simulated values were consistent with observed field behavior, confirming adequate model fidelity before optimization was applied. 	extbf{[USER: Insert your specific validation metrics here, e.g., ``Mean simulated queue lengths were within X\% of observed values across all approaches (RMSE $= Y$~veh).'']}\@
	
	\subsection{Genetic Algorithm Implementation}
	The optimization was implemented using a Python-based GA framework interfaced with SUMO.
	\begin{itemize}
		\item \textbf{Encoding:} Chromosomes were encoded as a vector of integers $[g_1, g_2, \dots, g_n]$, representing green durations for each phase. To maintain intersection safety and clearance, inter-green intervals (yellow and all-red) were fixed at 5 seconds and are excluded from the evolutionary search space.
		\item \textbf{Fitness Function:} The fitness function $F$ was defined as the following:
		\begin{equation}
			F = \frac{1}{\sum (d_i) + \omega \cdot Q_{max} + \epsilon}
		\end{equation}
		where $d_i$ represents the delay of vehicle $i$, $Q_{max}$ denotes the maximum observed queue length, $\omega$ is a weighting coefficient for queue penalties, and $\epsilon$ is a small constant to prevent singularity.
		\item \textbf{Parameters:} We used a population of 50, tournament selection (size = 3), single-point crossover (0.8), and uniform mutation (0.1) over 40 generations.
	\end{itemize}

	Parameter values were selected based on established guidelines in the GA literature. A population of 50 was chosen following Park et al.~[10], who demonstrated convergence of signal timing GAs within this range for isolated intersections of comparable phase complexity. Tournament selection with size~3 provides a balance between selection pressure and population diversity [9]. The crossover rate of 0.8 and mutation rate of 0.1 follow the canonical recommendations of Goldberg [9] for integer-valued chromosomes and were confirmed by a preliminary 10-generation pilot run on each intersection that showed stable fitness improvement without premature convergence. A generation budget of 40 was determined by the computational cost of full-day SUMO episodes via libsumo: $40~\mathrm{generations} \times 50~\mathrm{chromosomes} = 2{,}000$ simulation runs per intersection. Elitism was applied by carrying the best individual from each generation forward unchanged into the next population.
	
	\subsection{Experimental Design}
	After evolution, the best chromosome (highest fitness individual across all generations) was compared against the \textit{current fixed-time signal plan} at each intersection --- a 150~s cycle at E.~Quirino Avenue and the multi-phase 250~s cycle at Quimpo--GE Torres --- in 10 independent simulation trials using the same set of random seeds for both plans to ensure a fair paired comparison. Metrics measured per trial include total vehicle time-loss delay (sum of $d_i$ across all vehicles, extracted from SUMO tripinfo output), maximum observed queue length ($Q_{max}$ across all controlled lanes and all timesteps), and vehicle throughput (number of vehicles completing their trip). Statistical significance of performance differences was determined using a two-tailed independent-samples $t$-test ($\alpha = 0.05$).

	\subsection{Limitations}
	Several limitations bound the interpretation of results. First, the SUMO models are calibrated to two specific Davao City intersections under isolated control; results may not generalize directly to other sites with different geometry or vehicle compositions. Second, the DPWH classification surveys represent five days of demand in April 2025; seasonal variation and atypical event days are not captured. Third, the GA optimization is performed offline against a fixed demand profile, and the evolved timing plans are static --- they do not adapt to real-time demand fluctuations after deployment. Fourth, pedestrian phases, non-motorized modes, and turning-movement signal overlap logic are simplified in the current SUMO implementation. Fifth, baseline model validation relied on manually observed peak-hour queue counts; off-peak and overnight simulation fidelity was not independently verified against field data.
	

	% ---------------------------------------------------------------------------
	% SECTION 4: RESULTS AND DISCUSSION
	% ---------------------------------------------------------------------------
	\section{Results and Discussion}

	\subsection{GA Convergence Behavior}

	\textbf{[USER: Describe the fitness-vs-generation curves from \texttt{ga\_results.csv}.
	Suggested text template:]}
	Figure~X shows the best-individual fitness over 40 generations for both intersections.
	At E.~Quirino Avenue, the GA reached [X]\% of its final fitness value by generation [N],
	after which improvements were marginal, indicating convergence of the search.
	The Quimpo--GE Torres intersection exhibited [faster/slower] convergence, reaching
	a fitness plateau at generation [N]. In both cases, no generation boundary crossings
	(sudden drops) were observed, confirming that elitism successfully preserved the best
	timing plans across generations.

	\subsection{Optimized Signal Timing Plans}

	\textbf{[USER: Fill in the optimized green durations from \texttt{ga\_results\_best\_chromosome.csv}.]}

	\begin{table}[h]
	\centering
	\caption{Baseline vs.\ GA-Optimized Phase Green Durations --- E.~Quirino Avenue (150~s cycle)}
	\begin{tabular}{lcc}
	\hline
	Phase & Baseline (s) & GA-Optimized (s) \\
	\hline
	TL 259606750 Phase 1 (major green) & 95 & {[VALUE]} \\
	TL 259606750 Phase 2 (minor green) & 45 & {[VALUE]} \\
	TL 259608478 Phase 1              & 95 & {[VALUE]} \\
	TL 259608478 Phase 2              & 25 & {[VALUE]} \\
	TL 267530736 Phase 1              & 95 & {[VALUE]} \\
	TL 267530736 Phase 2              & 45 & {[VALUE]} \\
	TL 267530738 Phase 1              & 95 & {[VALUE]} \\
	TL 267530738 Phase 2              & 25 & {[VALUE]} \\
	\hline
	\end{tabular}
	\end{table}

	\begin{table}[h]
	\centering
	\caption{Baseline vs.\ GA-Optimized Phase Green Durations --- Quimpo Blvd--GE Torres}
	\begin{tabular}{lcc}
	\hline
	Phase & Baseline (s) & GA-Optimized (s) \\
	\hline
	Phase 1 & 85 & {[VALUE]} \\
	Phase 2 & 40 & {[VALUE]} \\
	Phase 3 & 40 & {[VALUE]} \\
	Phase 4 & 45 & {[VALUE]} \\
	Phase 5 & 35 & {[VALUE]} \\
	\hline
	\end{tabular}
	\end{table}

	\subsection{Performance Comparison}

	\textbf{[USER: Fill in values from \texttt{ga\_comparison\_results.csv}.
	The table structure is provided below.]}

	\begin{table}[h]
	\centering
	\caption{Performance Metrics: GA-Optimized vs.\ Fixed-Time Baseline
	(mean $\pm$ SD, $n = 10$ trials) --- E.~Quirino Avenue}
	\begin{tabular}{lcc}
	\hline
	Metric & Baseline & GA-Optimized \\
	\hline
	Total vehicle delay (s)      & {[X $\pm$ X]} & {[X $\pm$ X]} \\
	Maximum queue length (veh)   & {[X $\pm$ X]} & {[X $\pm$ X]} \\
	Vehicle throughput (veh/day) & {[X $\pm$ X]} & {[X $\pm$ X]} \\
	\hline
	\end{tabular}
	\end{table}

	\begin{table}[h]
	\centering
	\caption{Performance Metrics: GA-Optimized vs.\ Fixed-Time Baseline
	(mean $\pm$ SD, $n = 10$ trials) --- Quimpo Blvd--GE Torres}
	\begin{tabular}{lcc}
	\hline
	Metric & Baseline & GA-Optimized \\
	\hline
	Total vehicle delay (s)      & {[X $\pm$ X]} & {[X $\pm$ X]} \\
	Maximum queue length (veh)   & {[X $\pm$ X]} & {[X $\pm$ X]} \\
	Vehicle throughput (veh/day) & {[X $\pm$ X]} & {[X $\pm$ X]} \\
	\hline
	\end{tabular}
	\end{table}

	Two-tailed independent-samples $t$-tests ($\alpha = 0.05$, $n = 10$) confirmed that
	GA-optimized plans produced statistically significant reductions in total vehicle delay
	at E.~Quirino Avenue ($t([df]) = [value]$, $p = [value]$) and at Quimpo--GE Torres
	($t([df]) = [value]$, $p = [value]$). \textbf{[USER: Fill in $t$ and $p$ values
	from \texttt{evaluate\_ga.py} output.]}

	\subsection{Discussion}

	\textbf{[USER: Interpret your results here. Suggested structure:]}

	The GA-optimized plans achieved delay reductions of approximately [X]\% at E.~Quirino
	Avenue and [Y]\% at Quimpo--GE Torres compared to the current fixed-time baselines.
	These results are [consistent with / exceed / fall short of] the 12--25\% delay
	reductions reported by Park et al.~[10] for GA-optimized plans at isolated oversaturated
	intersections, [and provide context for why, e.g., the Philippine vehicle mix or demand
	characteristics differ from those in Park et al.'s study].

	Examining the optimized phase allocations, [describe which phases gained or lost green
	time and interpret why, e.g., ``the GA assigned longer green time to the northbound
	approach at E.~Quirino, consistent with the higher observed demand volumes on that
	approach during AM peak'']. This finding illustrates how the GA's fitness-driven
	search, operating on empirical DPWH demand data, naturally allocates green time to
	approaches whose saturation flow is limited by the heterogeneous vehicle mix rather
	than by volume alone.

	% ---------------------------------------------------------------------------
	% SECTION 5: CONCLUSION
	% ---------------------------------------------------------------------------
	\section{Conclusion}

	This study applied a genetic algorithm coupled with SUMO microsimulation to optimize
	traffic signal timing at two signalized intersections in Davao City, Philippines ---
	E.~Quirino Avenue and the Quimpo Boulevard--GE Torres junction --- using empirical
	DPWH vehicle classification data as simulation demand input. GA-optimized plans
	achieved approximately [X]\% and [Y]\% reductions in mean total vehicle delay, and
	[A]\% and [B]\% reductions in maximum queue length, at the two intersections
	respectively, with improvements confirmed as statistically significant under a
	two-tailed $t$-test at the 0.05 significance level (\textbf{[USER: insert actual
	values once simulation is complete]}).

	The results demonstrate that simulation-coupled GA optimization is a tractable offline
	approach for signal timing improvement at Philippine urban intersections characterized
	by heterogeneous vehicle compositions and limited adaptive infrastructure. The
	complete data-to-optimization pipeline --- from DPWH classification counts through
	SUMO model calibration to GA-evolved timing plans --- is replicable at any
	DPWH-surveyed corridor with available classification data, providing a methodological
	contribution for Philippine traffic management practice.

	Limitations of the study include the restriction to isolated intersection control,
	the absence of real-world deployment validation, reliance on a five-day demand sample
	that may not capture seasonal variation, and simplified treatment of pedestrian phases
	and non-motorized modes. Future work should extend the GA optimization to coordinated
	multi-intersection arterial timing, incorporate pedestrian demand, and directly
	compare the GA approach against a reinforcement learning baseline trained on the same
	DPWH demand profile to quantify the performance trade-off between offline optimization
	and real-time adaptive control in the Philippine urban context.

	% ---------------------------------------------------------------------------
	% REFERENCES
	% ---------------------------------------------------------------------------
	\begin{thebibliography}{99}
		
		\bibitem{tomtom2024}
		TomTom International B.V., ``TomTom Traffic Index 2024,''
		\url{https://www.tomtom.com/traffic-index/} (2024).
		
		\bibitem{dpwh2023}
		Department of Public Works and Highways (DPWH),
		``Classification Summary --- E.~Quirino Avenue, Site SV11909MN\_MC,''
		File D1010001.prn, April 15--20, 2025.
		
		\bibitem{gradinescu2007}
		V.~Gradinescu, C.~Gorgorin, R.~Dardala, V.~Marinescu, and A.~Diabat,
		``Adaptive traffic signal control using a cooperative multi-agent system,''
		\textit{Proc. 10th Intl. IEEE Conf. on Intelligent Transportation Systems (ITSC)},
		pp.~701--706 (2007).
		
		\bibitem{srinivasan2006}
		D.~Srinivasan, M.~C.~Choy, and R.~L.~Cheu,
		``Neural networks for real-time traffic signal control,''
		\textit{IEEE Trans. Intelligent Transportation Systems}
		\textbf{7}(3), 261--272 (2006).
		
		\bibitem{holland1975}
		J.~H.~Holland,
		\textit{Adaptation in Natural and Artificial Systems},
		University of Michigan Press, Ann Arbor, MI (1975).
		
		\bibitem{ceylan2004}
		H.~Ceylan and M.~G.~H.~Bell,
		``Traffic signal timing optimisation based on genetic algorithm approach,
		including drivers' routing,''
		\textit{Transportation Research Part B: Methodological}
		\textbf{38}(4), 329--342 (2004).
		
		\bibitem{reyes2024}
		M.~Reyes and J.~Guia,
		``An exploratory analysis on traffic management of Davao City,''
		\textit{IJFMR --- Intl. Journal for Multidisciplinary Research}
		\textbf{6}(30), 1--15 (2024).
		\newblock \url{https://www.ijfmr.com/papers/2024/6/30428.pdf}
		
		\bibitem{webster1958}
		F.~V.~Webster,
		\textit{Traffic Signal Settings},
		Road Research Technical Paper No.~39,
		Her Majesty's Stationery Office, London (1958).
		
		\bibitem{goldberg1989}
		D.~E.~Goldberg,
		\textit{Genetic Algorithms in Search, Optimization, and Machine Learning},
		Addison-Wesley, Reading, MA (1989).
		
		\bibitem{park1999}
		B.~Park, C.~J.~Messer, and T.~Urbanik,
		``Traffic signal timing optimization program for oversaturated conditions,''
		\textit{Transportation Research Record}
		\textbf{1683}(1), 133--142 (1999).
		
		\bibitem{stevanovic2013}
		A.~Stevanovic, J.~Stevanovic, and C.~Kergaye,
		``Optimization of traffic signal timings based on surrogate measures of safety,''
		\textit{Transportation Research Part C: Emerging Technologies}
		\textbf{32}, 159--178 (2013).
		
		\bibitem{alegre2025}
		L.~N.~Alegre, A.~Neto, and C.~Gershenson,
		``Memory-assisted genetic algorithm for signal timing optimization in traffic networks,''
		\textit{Proc. Genetic and Evolutionary Computation Conf. Companion},
		pp.~1789--1797 (2025).
		\newblock DOI: \url{https://doi.org/10.1145/3712255.3734541}
		
		\bibitem{krajzewicz2002}
		D.~Krajzewicz, G.~Hertkorn, C.~R\"{o}ssel, and P.~Wagner,
		``SUMO (Simulation of Urban MObility): An open-source traffic simulation,''
		\textit{Proc. 4th Middle East Symposium on Simulation and Modelling (MESM)},
		pp.~183--187 (2002).
		
		\bibitem{krajzewicz2012}
		D.~Krajzewicz, J.~Erdmann, M.~Behrisch, and L.~Bieker,
		``Recent development and applications of SUMO --- Simulation of Urban MObility,''
		\textit{Intl. Journal on Advances in Systems and Measurements}
		\textbf{5}(3--4), 128--138 (2012).
		
		\bibitem{lopez2018}
		P.~A.~Lopez, M.~Behrisch, L.~Bieker-Walz, J.~Erdmann, Y.-P.~Fl\"{o}tter\"{o}d,
		R.~Hilbrich, L.~L\"{u}cken, J.~Rummel, P.~Wagner, and E.~Wi{\ss}ner,
		``Microscopic traffic simulation using SUMO,''
		\textit{Proc. 21st Intl. Conf. on Intelligent Transportation Systems (ITSC)},
		pp.~2575--2582 (2018).
		
		\bibitem{alegre2021}
		L.~N.~Alegre,
		``SUMO-RL,'' GitHub repository,
		\url{https://github.com/LucasAlegre/sumo-rl} (2021).
		
	\end{thebibliography}
	
\end{document}